% Auto-generated by src/make_tables.py -- do not edit by hand.
\begin{table}[t]
\centering
\caption{The order of the reconstruction governs control accuracy, not the accuracy of the solution. Each row turns nodal values into a gradient a different way: a piecewise-linear (order-one) reconstruction whose gradient is probed by finite differences, that same reconstruction differentiated exactly, and a cubic (order-three) reconstruction probed exactly as the first. The columns report the convergence rate of the \emph{gradient} --- the control --- measured on functions whose gradients are known exactly, so the rates are properties of the reconstruction rather than of the benchmark; the first four columns are this paper's own problem, averaged over $p = 2, 3$ and $5$. Orders agreed across dimensions to within 2.0. The last column is the Merton value function, a power function from a genuinely non-linear control problem --- non-polynomial, so the cubic reconstruction is not exact there. Differentiating the piecewise-linear reconstruction exactly is no better than probing it, so consistency with the interpolant buys nothing. The quadratic column is this paper's own case and flatters the cubic method, which reproduces quadratics exactly; every other column is the one to quote. Errors are normalised by the spread of the true gradient.}
\label{tab:extract}
\small
\begin{tabular}{lrrrrr}
\toprule
 & \multicolumn{5}{c}{gradient-error convergence rate} \\
\cmidrule(lr){2-6}
reconstruction & quad. & quartic & exp. & trig. & Merton \\
\midrule
multilinear $+$ probe & 2.2 & 2.3 & 2.2 & 2.2 & 2.1 \\
multilinear, slope & 2.0 & 2.0 & 2.0 & 2.0 & 1.9 \\
cubic $+$ probe & exact & 4.8 & 4.4 & 5.6 & 5.2 \\
\bottomrule
\end{tabular}
\end{table}
